\chapter{L'idée JEPA : deviner l'idée, pas les pixels}
\label{chap:idee-jepa}

\begin{quote}\itshape
Nous avons laissé une question en suspens à la fin de la partie~I. Quand la
machine cache un morceau de ce qu'elle regarde et essaie de le deviner,
que doit-elle deviner au juste~? Les pixels qui manquent, un par un~? Ou
bien autre chose~? Ce chapitre répond à cette question, et la réponse est
le cœur de tout le reste du livre. Nous allons monter, pièce par pièce, la
machine qui porte cette réponse. Elle s'appelle JEPA.
\end{quote}

\begin{prerequis}
Les quatre chapitres de la partie~I, et rien d'autre. Il faut savoir qu'une
image est une grille de nombres (chapitre~\ref{chap:image-video})~; qu'un
modèle possède des boutons que l'on tourne pour faire baisser une erreur
(chapitre~\ref{chap:apprendre})~; qu'on peut apprendre sans étiquettes en
cachant un morceau de l'image, et qu'un encodeur transforme une image en un
petit paquet de nombres appelé représentation
(chapitre~\ref{chap:auto-supervision})~; et enfin comment une image devient
une file de jetons qui discutent entre eux
(chapitre~\ref{chap:vit}). Les seules opérations employées ici seront
l'addition, la soustraction, la multiplication et la division.
\end{prerequis}

Au chapitre~\ref{chap:auto-supervision}, nous avons trouvé un moyen
d'apprendre sans que personne n'écrive la bonne réponse~: on cache un
morceau de l'image, et on demande à la machine de deviner ce qui manque.
C'était une belle idée, et elle marchait déjà sur notre petit monde de huit
cases sur huit.

Puis, à la fin de ce même chapitre, nous avons buté sur une question que
nous avons laissée ouverte~: fallait-il deviner les \emph{pixels} cachés, ou
seulement \emph{l'idée} de ce qui manquait~? Nous avions même donné un avant-goût
de la réponse, et lâché un nom~: JEPA. Le moment est venu de tenir cette
promesse. Nous allons reprendre cette question, la trancher pour de bon, et
surtout monter la machine qui en découle.

\section{Deviner les pixels, ou deviner l'idée}

\begin{enclair}
Quand on cache un morceau d'image, on peut demander deux choses très
différentes à la machine. Soit on lui demande de \emph{redessiner} le
morceau manquant, pixel par pixel, comme un peintre qui refait un bout de
tableau. Soit on lui demande seulement de \emph{décrire} ce qu'il y avait,
comme quelqu'un qui dirait «~là, il y avait de l'herbe~». La première
demande est terriblement difficile. La seconde est raisonnable. V-JEPA
choisit la seconde.
\end{enclair}

Reprenons la pelouse du chapitre~\ref{chap:auto-supervision}. On a une photo
de jardin, on cache un carré au milieu, et ce carré tombe sur l'herbe.
Posons cette fois les deux questions côte à côte, bien en face l'une de
l'autre.

\textbf{Première question, celle du peintre.} «~Redessine-moi ce carré,
exactement comme il était.~» Nous savons déjà où cela mène~: il faudrait
retrouver chaque brin d'herbe, celui qui penche à gauche, celui qui est un
peu jauni, la petite ombre entre les deux. Ces détails-là ne se déduisent de
rien, ils sont le fruit du hasard, et la machine s'épuiserait à deviner ce
qui ne peut pas se deviner.

\textbf{Deuxième question, celle du témoin.} «~Dis-moi seulement ce qu'il y
avait dans ce carré.~» Cette fois, la machine peut répondre~: «~de l'herbe, verte, en
pleine lumière, et c'est la suite de la pelouse d'à côté~». Cette réponse-là,
elle peut la trouver. Elle est même obligée de comprendre la scène pour la
trouver. Et surtout, elle a le droit de \emph{ne pas savoir} où était
exactement chaque brin d'herbe, parce qu'on ne le lui demande pas.

\begin{intuition}
Imaginez qu'un ami vous téléphone depuis une pièce que vous ne voyez pas, et
qu'il vous demande de deviner ce qu'il y a derrière lui. Si vous devez dire
«~il y a une bibliothèque, avec des livres~», vous avez une chance de tomber
juste. Si vous devez dire le titre exact de chaque livre, dans l'ordre, vous
n'avez aucune chance. Le premier jeu vous apprend à réfléchir. Le second
ne vous apprend rien~: il vous fait seulement perdre à tous les coups.
\end{intuition}

\begin{definition}[Méthode générative][generative]
Une méthode est dite \emph{générative} quand elle demande au modèle de
reconstruire l'entrée elle-même, c'est-à-dire de ressortir des pixels. Le
mot vient de «~générer~», qui veut dire «~fabriquer~». On compare alors les
pixels fabriqués aux vrais pixels.
\end{definition}

\begin{definition}[Méthode prédictive dans l'espace latent][predictif-latent]
Une méthode est dite \emph{prédictive dans l'espace latent} quand elle
demande au modèle de prédire, non pas les pixels, mais la
\emph{représentation} de la zone manquante — c'est-à-dire le petit paquet de
nombres qui la résume. On rappelle que l'\emph{espace latent} est simplement
le nom savant du monde des représentations~: l'endroit où vivent ces paquets
de nombres, une fois que l'encodeur a fait son travail
(chapitre~\ref{chap:auto-supervision}).
\end{definition}

Voilà les deux chemins nommés. Le chapitre~\ref{chap:auto-supervision} nous
avait mis sur la piste du second~; nous savons maintenant comment il
s'appelle, et la figure~\ref{fig:jepa-deux-voies} les met côte à côte pour
qu'on ne les confonde plus jamais.

\begin{schema}{Les deux façons de remplir un trou. En haut, la méthode
générative redessine les pixels, et doit donc inventer chaque brin d'herbe.
En bas, la méthode JEPA se contente de décrire ce qui manquait, avec un
petit paquet de nombres.}[jepa-deux-voies]

% ---------- Ligne du haut : la voie générative ----------
\node[font=\sffamily\bfseries\scriptsize, text=mbInk, anchor=west]
      at (0,2.05) {La façon du peintre : redessiner les pixels};

\fill[mbAzure!18] (0,0.8) rectangle (2.2,1.6);
\fill[mbGreen!30] (0,0)   rectangle (2.2,0.8);
\draw[mbSlate, line width=0.7pt] (0,0) rectangle (2.2,1.6);
\fill[mbSlate!25] (0.75,0.4) rectangle (1.45,1.1);
\draw[mbSlate, dashed, line width=0.7pt] (0.75,0.4) rectangle (1.45,1.1);
\node[font=\sffamily\scriptsize, text=mbSlate] at (1.1,-0.35) {la photo trouée};

\draw[fleche] (2.45,0.8) -- (3.65,0.8);
\node[etiquette] at (3.05,1.1) {on redessine};

\fill[mbAzure!18] (3.9,0.8) rectangle (6.1,1.6);
\fill[mbGreen!30] (3.9,0)   rectangle (6.1,0.8);
\draw[mbSlate, line width=0.7pt] (3.9,0) rectangle (6.1,1.6);
\fill[mbGreen!30] (4.65,0.4) rectangle (5.35,1.1);
\draw[mbGreen!60!black, line width=0.6pt] (4.65,0.4) rectangle (5.35,1.1);
\foreach \k in {0,1,...,6}{
  \draw[draw=mbGreen!70!black, line width=0.5pt]
        (4.72+\k*0.09, 0.45) -- (4.67+\k*0.09, 0.95);
}
\node[font=\sffamily\scriptsize, text=mbSlate] at (5,-0.35) {chaque brin inventé};

\node[annotation bulle, align=center, anchor=west] (bulhaut) at (7.15,0.8)
     {il faudrait deviner\\ce qui ne se devine pas};
\draw[trait annotation, line width=0.7pt, draw=mbAmber!80]
      (6.25,0.8) -- (bulhaut.west);

% ---------- Ligne du bas : la voie JEPA ----------
\begin{scope}[yshift=-3.3cm]
\node[font=\sffamily\bfseries\scriptsize, text=mbInk, anchor=west]
      at (0,2.05) {La façon de JEPA : décrire l'idée};

\fill[mbAzure!18] (0,0.8) rectangle (2.2,1.6);
\fill[mbGreen!30] (0,0)   rectangle (2.2,0.8);
\draw[mbSlate, line width=0.7pt] (0,0) rectangle (2.2,1.6);
\fill[mbSlate!25] (0.75,0.4) rectangle (1.45,1.1);
\draw[mbSlate, dashed, line width=0.7pt] (0.75,0.4) rectangle (1.45,1.1);
\node[font=\sffamily\scriptsize, text=mbSlate] at (1.1,-0.35) {la photo trouée};

% Le nuage est posé AVANT les flèches : celles-ci s'accrochent ensuite à
% ses ancrages (nuage.west, nuage.east) au lieu de viser des coordonnées
% devinées à la main. Sans cela, leurs pointes se retrouvent cachées sous
% le nuage, qui déborde largement de sa largeur annoncée.
\node[cloud, cloud puffs=13, cloud puff arc=110, aspect=1.9,
      draw=mbAmber!80!black, fill=mbAmber!12, line width=0.7pt,
      minimum width=2.7cm, minimum height=1.6cm,
      font=\sffamily\scriptsize, text=mbInk] (nuage) at (5.05,0.8)
     {$0{,}8$ \ $-0{,}3$ \ $1{,}2$};

\draw[fleche accent, shorten >=2pt] (2.45,0.8) -- (nuage.west);
\node[etiquette] at (3.05,1.1) {on décrit};

\node[font=\sffamily\scriptsize, text=mbSlate] at (5.05,-0.35)
     {un petit paquet de nombres};

\node[annotation bulle, align=center, anchor=west] (bulbas) at (7.15,0.8)
     {« de l'herbe verte,\\en pleine lumière » :\\c'est faisable};
\draw[trait annotation, line width=0.7pt, draw=mbAmber!80, shorten <=2pt]
      (nuage.east) -- (bulbas.west);
\end{scope}
\end{schema}

\begin{piege}
Une erreur revient tout le temps, alors autant l'écarter tout de suite~: on
croit souvent que la voie JEPA est meilleure parce qu'elle aurait
\emph{moins de nombres} à deviner. C'est faux, et il faut le savoir. Un
carré de $64 \times 64$ pixels en couleur, cela fait
$64 \times 64 \times 3 = 12\,288$ nombres à redessiner. La représentation de
ce même carré peut très bien en contenir davantage. Ce n'est donc pas une
question de \emph{quantité}. C'est une question de \emph{nature}~: les
pixels sont imposés de l'extérieur, avec leur part de hasard, alors que la
représentation est fabriquée par la machine elle-même, qui a donc le droit
de laisser de côté ce qui n'est pas devinable.
\end{piege}

\begin{aretenir}
Prédire des pixels oblige à deviner du hasard. Prédire une représentation
permet de l'oublier. C'est là toute la différence, et c'est elle qui donne
son nom à toute la famille JEPA.
\end{aretenir}

Nous savons maintenant \emph{quoi} deviner. Voyons à partir de quoi, et
mettons enfin des noms sur les deux images qui entrent dans la machine.

\section{Les deux vues : la trouée et la complète}

\begin{enclair}
La machine regarde toujours la même vidéo deux fois. La première fois, on
lui en cache des morceaux~: c'est la vue trouée. La deuxième fois, on la lui
montre en entier~: c'est la vue complète. La première sert de question, la
seconde sert de corrigé.
\end{enclair}

Reprenons le jeu du chapitre~\ref{chap:auto-supervision}, mais en donnant
cette fois un nom précis à chaque chose. On part d'une vidéo. On en fabrique
deux versions.

\begin{definition}[La vue trouée, notée $x$][vue-trouee]
La \emph{vue trouée} est la vidéo à laquelle on a retiré des morceaux. On la
note $x$. Les morceaux retirés s'appellent les zones \emph{masquées}, et les
morceaux que l'on a laissés s'appellent le \emph{contexte}. On dit aussi que
$x$ est la vue \emph{corrompue}, ce qui veut simplement dire «~abîmée
exprès~».
\end{definition}

\begin{definition}[La vue complète, notée $y$][vue-complete]
La \emph{vue complète} est la même vidéo, sans rien retirer du tout. On la
note $y$. On l'appelle aussi la vue \emph{propre}, au sens de «~intacte~».
\end{definition}

Attention à ne pas se tromper sur un point~: $x$ et $y$ ne sont pas deux
vidéos différentes. C'est la même vidéo, la même seconde, la même scène. La
seule différence, c'est qu'on a caché des morceaux dans l'une et pas dans
l'autre.

\begin{exemple}[Compter les jetons cachés]
Reprenons un clip de la phase principale d'entraînement, celui du
chapitre~\ref{chap:vit}~: seize images en $256 \times 256$, avec des patchs
de $16$ et des tubelets de $2$. Nous avions compté
$8 \times 16 \times 16 = 2\,048$ jetons.

Disons que le masque en recouvre $1\,200$. Alors~:

\begin{center}
\footnotesize
\begin{tabular}{lr}
\toprule
\textbf{Ce que l'on compte} & \textbf{Combien} \\
\midrule
Jetons en tout dans le clip        & $2\,048$ \\
Jetons masqués (le trou)           & $1\,200$ \\
Jetons de contexte (ce qui reste)  & $848$ \\
\bottomrule
\end{tabular}
\end{center}

La machine ne verra donc que $848$ jetons, et devra parler des $1\,200$
autres. Retenez ces nombres~: ils vont nous resservir tout au long du
chapitre.
\end{exemple}

\begin{schema}{La même vidéo, vue deux fois. À gauche la vue trouée $x$,
où seuls les jetons bleus sont visibles. À droite la vue complète $y$, où
tout est visible.}[jepa-vues]

% --- vue trouée ---
\foreach \i in {0,1,2,3,4,5}{
  \foreach \j in {0,1,2,3,4}{
    \fill[mbAzure!25, draw=mbAzure!60, line width=0.4pt]
          (\i*0.5, \j*0.5) rectangle ++(0.46,0.46);
  }
}
\foreach \i in {1,2,3}{
  \foreach \j in {1,2,3}{
    \fill[mbSlate!30, draw=mbSlate!70, line width=0.4pt, dashed]
          (\i*0.5, \j*0.5) rectangle ++(0.46,0.46);
  }
}
\node[font=\sffamily\small, text=mbInk] at (1.5,-0.55)
     {la vue trouée $x$};
\node[etiquette, text=mbAzure!80!black] at (1.5,-1.05) {bleu : le contexte};
\node[etiquette, text=mbSlate] at (1.5,-1.45) {gris : les zones masquées};

% --- vue complète ---
\begin{scope}[xshift=5.5cm]
\foreach \i in {0,1,2,3,4,5}{
  \foreach \j in {0,1,2,3,4}{
    \fill[mbTeal!22, draw=mbTeal!60, line width=0.4pt]
          (\i*0.5, \j*0.5) rectangle ++(0.46,0.46);
  }
}
\node[font=\sffamily\small, text=mbInk] at (1.5,-0.55)
     {la vue complète $y$};
\node[etiquette, text=mbTeal!80!black] at (1.5,-1.05) {tout est visible};
\end{scope}
\end{schema}

Nous avons nos deux vues. Il faut maintenant deux appareils pour les
regarder. Et là, une petite surprise nous attend~: ces deux appareils sont
bâtis de la même façon, mais leurs boutons ne sont pas réglés pareil.

\section{L'élève et le professeur : deux encodeurs, presque jumeaux}

\begin{enclair}
La vue trouée est regardée par un premier encodeur, que nous appellerons
l'\emph{élève}. La vue complète est regardée par un second encodeur, que
nous appellerons le \emph{professeur}. L'élève est celui qui apprend~: on
lui tourne ses boutons à chaque pas. Le professeur, lui, ne travaille pas~:
il se contente de recopier l'élève, très lentement.
\end{enclair}

Avant d'aller plus loin, il faut une convention d'écriture. Nous avons vu au
chapitre~\ref{chap:apprendre} qu'un modèle est une machine à boutons, et
qu'apprendre revient à tourner ces boutons. Un gros modèle en a des
millions. Plutôt que de les nommer un par un, on leur donne un seul nom
collectif, avec une lettre grecque.

\begin{notation}[Les lettres grecques des boutons][lettres-grecques]
\begin{itemize}
  \item $\theta$ se lit «~thêta~». C'est le nom collectif de \emph{tous} les
        boutons de l'encodeur élève. Écrire $E_{\theta}$ veut donc dire
        «~l'encodeur, réglé avec les boutons $\theta$~».
  \item $\bar{\theta}$ se lit «~thêta barre~». C'est le nom collectif de
        tous les boutons de l'encodeur professeur. La petite barre au-dessus
        rappelle qu'il s'agit d'une \emph{moyenne} de l'élève~: nous verrons
        cela en détail dans la section~\ref{sec:ema}.
  \item $\varphi$ se lit «~phi~». C'est le nom collectif des boutons du
        prédicteur, la pièce que nous montons juste après.
\end{itemize}
Aucune de ces lettres ne cache une difficulté. Ce sont des étiquettes, rien
de plus. On aurait pu les appeler $A$, $B$ et $C$.
\end{notation}

Voici ce que fait chacun des deux encodeurs.

\textbf{L'encodeur élève, $E_{\theta}$.} Il reçoit la vue trouée $x$,
c'est-à-dire seulement les $848$ jetons de contexte de notre exemple. Il ne
voit rien d'autre. Il les fait discuter entre eux, comme nous l'avons vu au
chapitre~\ref{chap:vit}, et il ressort $848$ jetons enrichis. C'est lui, et
lui seul, qui apprend. À chaque pas d'entraînement, ses boutons bougent.

\textbf{L'encodeur professeur, $E_{\bar{\theta}}$.} Il reçoit la vue
complète $y$, donc les $2\,048$ jetons, sans aucun trou. Il ressort $2\,048$
jetons enrichis. Ces jetons-là serviront de \emph{corrigé}~: ce sont les
réponses que l'élève doit essayer de retrouver. Le professeur, lui,
n'apprend pas au sens habituel~: aucun calcul d'erreur ne vient tourner ses
boutons. Il se contente de suivre l'élève de très loin.

\begin{intuition}
Pensez à un élève qui fait un devoir avec la moitié des pages arrachées, et
à un professeur qui a le sujet complet sous les yeux. Le professeur ne fait
pas d'effort~: il lit. Mais c'est sa copie qui sert de référence. Et le
professeur n'est pas quelqu'un d'autre~: c'est le même élève, quelques
milliers de pas plus tôt, une version un peu plus posée de lui-même.
\end{intuition}

\begin{notation}[Les autres noms de ces deux encodeurs][x-encodeur]
Dans l'article, et dans le code publié sur internet, ces deux encodeurs
portent des noms tirés des lettres des deux vues. L'encodeur élève, celui
qui regarde la vue trouée $x$, s'appelle le \emph{$x$-encodeur}. L'encodeur
professeur, celui qui regarde la vue complète $y$, s'appelle le
\emph{$y$-encodeur}. Ce sont exactement les mêmes pièces que notre élève et
notre professeur~: seul le nom change. Nous emploierons les deux façons de
dire, pour que vous soyez à l'aise le jour où vous ouvrirez le code.
\end{notation}

\begin{schema}{La paire d'encodeurs. Les deux appareils sont bâtis de la
même façon, mais un seul des deux apprend~: le professeur se contente de
recopier l'élève, très lentement.}[jepa-paire]

\node[bloc, font=\sffamily\scriptsize, minimum width=28mm] (el) at (0,0)
     {encodeur élève\\$E_{\theta}$\\[1pt]\scriptsize il apprend};
\node[bloc calcul, font=\sffamily\scriptsize, minimum width=28mm]
     (pr) at (0,-2.6)
     {encodeur professeur\\$E_{\bar{\theta}}$\\[1pt]\scriptsize il recopie};

\node[donnee, font=\sffamily\scriptsize, minimum width=15mm]
     (vx) at (-3.7,0) {vue trouée\\$x$};
\node[donnee, font=\sffamily\scriptsize, minimum width=15mm]
     (vy) at (-3.7,-2.6) {vue complète\\$y$};
\draw[fleche] (vx) -- (el);
\draw[fleche] (vy) -- (pr);

\draw[fleche] (el.east) -- ++(1.0,0)
      node[right, font=\sffamily\scriptsize, text=mbInk, align=left]
      {848 jetons\\enrichis};
\draw[fleche] (pr.east) -- ++(1.0,0)
      node[right, font=\sffamily\scriptsize, text=mbInk, align=left]
      {2 048 cibles};

\draw[fleche pointillee] (el) -- (pr);
\node[etiquette, align=center] at (-1.35,-1.3)
     {recopie\\très lentement};
\end{schema}

\begin{definition}[Cible][cible]
On appelle \emph{cible} la sortie du professeur~: c'est le paquet de nombres
que la machine essaie de retrouver. Le mot est celui du tir à l'arc~: c'est
ce que l'on vise. Attention à ne pas la confondre avec les \emph{étiquettes}
du chapitre~\ref{chap:auto-supervision}, qui étaient écrites par des
humains~: ici, la cible est fabriquée par la machine elle-même, et personne
ne l'a jamais écrite nulle part.
\end{definition}

Une question doit vous venir à l'esprit, et elle est excellente~: pourquoi
faire deux encodeurs, alors qu'un seul suffirait~? Après tout, on pourrait
donner $x$ et $y$ au même appareil. Gardez cette question bien au chaud~: la
réponse arrive à la section~\ref{sec:effondrement}, et c'est une des plus
jolies idées du livre.

Pour l'instant, il nous manque encore une pièce, et c'est celle qui fait le
lien entre les deux côtés.

\section{Le prédicteur : la pièce qui fait le pont}

\begin{enclair}
L'élève ne voit que les morceaux visibles. Il ne peut donc rien dire des
morceaux cachés. Il faut une pièce supplémentaire, qui prenne ce que l'élève
a compris et qui en tire une réponse pour les endroits qu'il n'a jamais vus.
Cette pièce s'appelle le \emph{prédicteur}.
\end{enclair}

Cette pièce, nous l'avons déjà croisée au
chapitre~\ref{chap:auto-supervision}~: c'est elle qui, dans notre petit
programme, partait du résumé et tentait de deviner la moitié cachée de
l'image. Nous la retrouvons ici, en beaucoup plus grand. Il est donc temps
de lui donner son écriture officielle.

\begin{notation}[Le prédicteur s'écrit $P_{\varphi}$][notation-predicteur]
On note le prédicteur $P_{\varphi}$, ce qui se lit «~$P$ de phi~» et veut
dire~: le prédicteur, réglé avec les boutons $\varphi$. Le
chapitre~\ref{chap:auto-supervision} nous avait déjà prévenus d'une chose
qu'il ne faut pas oublier~: à la fin de l'entraînement, c'est l'encodeur que
l'on garde précieusement, pas le prédicteur. Une chose change tout de même
ici~: notre prédicteur ne rend plus une seule réponse, mais une réponse par
jeton reçu.
\end{notation}

Le prédicteur reçoit deux choses, et il faut bien les distinguer.

\textbf{Un~: les jetons de contexte}, c'est-à-dire les $848$ jetons que
l'élève vient d'enrichir. Ce sont les indices dont on dispose.

\textbf{Deux~: les jetons de masque.} Voilà une nouveauté qu'il faut
expliquer, car elle surprend toujours. Pour les $1\,200$ endroits cachés, on
ne peut évidemment pas donner le contenu, puisque c'est justement ce que
l'on cache. On donne alors à la place un jeton spécial, qui ne dit qu'une
seule chose~: «~ici, à cet endroit précis, il manquait quelque chose~».

\begin{definition}[Jeton de masque, noté $\Delta_y$][jeton-masque]
Un \emph{jeton de masque} est un jeton qui ne contient aucune information
sur l'image, mais qui porte une adresse~: la position spatiale et temporelle
du morceau caché. Son contenu est un paquet de nombres appris pendant
l'entraînement, le même pour tous les trous~; seule l'adresse change. On
note l'ensemble de ces jetons $\Delta_y$, et $\Delta$ se lit «~delta~».
\end{definition}

\begin{intuition}
Un jeton de masque, c'est une enveloppe vide sur laquelle on a écrit
l'adresse. On la tend au prédicteur en disant~: «~remplis-la~». Sans
l'adresse, le prédicteur ne saurait pas de quel trou on parle, et il
répondrait la même chose pour tous.
\end{intuition}

Le prédicteur met bout à bout ces deux paquets de jetons, les fait discuter
entre eux — c'est encore un transformeur, exactement comme au
chapitre~\ref{chap:vit} — et ressort un jeton pour chaque jeton d'entrée.

\begin{exemple}[Le compte exact des jetons dans le prédicteur]
Avec nos chiffres, le prédicteur reçoit~:

\[
  \underbrace{848}_{\text{jetons de contexte}}
  \;+\;
  \underbrace{1\,200}_{\text{jetons de masque}}
  \;=\;
  2\,048 \ \text{jetons en entrée},
\]

et il ressort $2\,048$ jetons en sortie. Notez bien ce résultat, car il va
compter énormément dans quelques chapitres~: \textbf{le prédicteur produit
une réponse pour tout le monde, y compris pour les endroits qui étaient déjà
visibles.} Nous verrons au chapitre suivant que V-JEPA~2 jette une bonne
partie de ces réponses, et au chapitre~10 que c'est précisément là que se
cache le défaut que V-JEPA~2.1 vient corriger.
\end{exemple}

Nous avons enfin toutes les pièces. Assemblons-les.

\section{Le schéma complet, monté sous vos yeux}

Voici la machine entière, figure~\ref{fig:jepa-central}. C'est le dessin le
plus important du livre~: tous les chapitres suivants ne feront que le
compléter ou le corriger. Prenez le temps de le suivre du doigt, de gauche à
droite.

\begin{schema}{L'architecture JEPA au complet. En haut, le chemin de
l'élève, qui part de la vue trouée. En bas, celui du professeur, qui part de
la vue complète. À droite, on compare les deux~: cet écart est la perte. La
barre rouge marque l'endroit où la correction n'a pas le droit de
passer.}[jepa-central]

% ---- chemin de l'élève (en haut) ----
\node[donnee, font=\sffamily\scriptsize, minimum width=16mm] (x) at (0,0)
     {vue trouée\\$x$};
\node[bloc, font=\sffamily\scriptsize, minimum width=19mm] (enc) at (2.9,0)
     {encodeur élève\\$E_{\theta}$};
\node[bloc accent, font=\sffamily\scriptsize, minimum width=19mm]
     (pred) at (6.0,0) {prédicteur\\$P_{\varphi}$};
\node[bloc, font=\sffamily\scriptsize, minimum width=16mm, draw=mbAmber,
      fill=mbAmber!8] (sortie) at (8.9,0) {prédiction};

\draw[fleche] (x) -- (enc);
\draw[fleche] (enc) -- (pred);
\node[etiquette] at (4.45,0.42) {848 jetons};
\draw[fleche] (pred) -- (sortie);

\node[bloc, font=\sffamily\scriptsize, minimum width=19mm, minimum height=6mm,
      draw=mbSlate!70, fill=mbMist] (masque) at (6.0,1.7)
     {jetons de masque $\Delta_y$};
\draw[fleche] (masque) -- (pred);

% ---- chemin du professeur (en bas) ----
\node[donnee, font=\sffamily\scriptsize, minimum width=16mm] (y) at (0,-3.3)
     {vue complète\\$y$};
\node[bloc calcul, font=\sffamily\scriptsize, minimum width=19mm]
     (encp) at (2.9,-3.3) {encodeur professeur\\$E_{\bar{\theta}}$};
\node[bloc calcul, font=\sffamily\scriptsize, minimum width=16mm]
     (target) at (8.9,-3.3) {cible};

\draw[fleche] (y) -- (encp);
\draw[fleche] (encp) -- (target);
\node[etiquette] at (5.4,-2.88) {2 048 jetons};

% ---- la comparaison ----
\draw[<->, draw=mbRose!85!black, line width=1pt] (sortie) -- (target);
\node[annotation bulle, align=center, anchor=east] at (8.35,-1.65)
     {on compare :\\cet écart,\\c'est la perte};

% ---- la copie lente ----
\draw[fleche pointillee] (enc) -- (encp);
\node[etiquette, align=center, anchor=east] at (2.72,-1.65)
     {copie très lente\\(moyenne mobile)};

% ---- le gradient bloqué ----
\draw[line width=2pt, draw=mbRose!85!black] (7.3,-3.75) -- (7.3,-2.85);
\node[annotation, text=mbRose!85!black, align=center] at (7.3,-4.35)
     {gradient bloqué :\\rien ne revient par ici};
\end{schema}

Récapitulons le trajet en cinq phrases, dans l'ordre du dessin.

\begin{enumerate}
  \item La vue trouée $x$ entre dans l'encodeur élève, qui en ressort les
        jetons de contexte enrichis.
  \item Ces jetons sont mis bout à bout avec les jetons de masque, qui
        portent les adresses des trous.
  \item Le prédicteur fait discuter tout ce beau monde et produit une
        prédiction pour chaque position.
  \item Pendant ce temps, la vue complète $y$ est passée dans l'encodeur
        professeur, qui produit les cibles.
  \item On compare la prédiction à la cible. L'écart est la \emph{perte}, et
        c'est lui que l'entraînement fait baisser.
\end{enumerate}

Il reste, sur ce dessin, deux choses que nous n'avons pas encore
expliquées~: la flèche en pointillé de la copie lente, et la barre rouge du
gradient bloqué. Elles ont l'air de petits détails de dessin. Ce sont en
réalité deux \emph{garde-fous}, et sans eux tout s'écroule.

Le mot mérite une explication, car nous allons l'employer sans arrêt. Un
garde-fou, dans la vie de tous les jours, c'est la barrière que l'on pose au
bord d'un pont pour empêcher les gens de tomber. Elle ne sert à rien tant
que tout va bien~; elle sert le jour où quelqu'un glisse. Nos deux
garde-fous sont exactement cela~: deux barrières qui ne gênent personne tant
que la machine travaille honnêtement, et qui l'empêchent de tomber le jour
où elle essaie de tricher. Nous allons voir pourquoi, mais donnons d'abord
son vrai nom à cette machine.

\section{JEPA : le nom décodé mot à mot}

\begin{enclair}
Cette architecture porte un nom très court, fait de quatre lettres~: JEPA.
Chacune de ces lettres est la première lettre d'un mot anglais. Nous allons
les traduire un par un, et vous verrez qu'il n'y a rien d'autre, derrière
ces quatre lettres, que la machine que nous venons de monter.
\end{enclair}

Ce nom s'écrit JEPA et se prononce «~jé-pa~». Il vient de \emph{Joint
Embedding Predictive Architecture}, une idée proposée
par~\textcite{lecun2022path}. Voici le décodage.

\begin{center}
\footnotesize
\begin{tabular}{llp{6.6cm}}
\toprule
\textbf{Lettre} & \textbf{Mot anglais} & \textbf{Ce que cela veut dire ici} \\
\midrule
J & \emph{Joint} &
«~conjoint~», c'est-à-dire à deux. Il y a deux chemins, celui de l'élève et
celui du professeur, et ils travaillent ensemble. \\[2pt]
E & \emph{Embedding} &
le \emph{plongement}, c'est-à-dire le fait de transformer une image en un
paquet de nombres. Nous l'avons vu au chapitre~\ref{chap:vit}. \\[2pt]
P & \emph{Predictive} &
«~prédictive~»~: il y a une pièce qui devine, le prédicteur. \\[2pt]
A & \emph{Architecture} &
la façon dont les pièces sont assemblées. Ni plus, ni moins. \\
\bottomrule
\end{tabular}
\end{center}

\begin{definition}[JEPA][jepa]
Une \emph{architecture prédictive à plongements conjoints}, ou JEPA, est un
assemblage où deux encodeurs transforment deux vues d'une même donnée en
paquets de nombres, et où un prédicteur essaie de retrouver le paquet du
second à partir du paquet du premier. La comparaison se fait toujours entre
paquets de nombres, jamais entre pixels.
\end{definition}

\begin{enia}
Cette idée n'est pas propre à la vidéo. On l'a appliquée au son, aux images
fixes avec I-JEPA~\parencite{assran2023ijepa}, puis à la vidéo avec la
famille V-JEPA, dont le «~V~» veut dire \emph{video}. V-JEPA~2.1, le modèle
de ce livre, est le dernier membre de cette famille. Toute la partie~IV
consistera à ajouter des pièces à ce squelette, sans jamais le changer.
\end{enia}

Le nom est posé, la machine est montée. Nous pouvons maintenant aborder le
danger qui menace tout ce montage — et comprendre enfin pourquoi il faut
deux encodeurs plutôt qu'un.

\section{Le grand danger : l'effondrement}
\label{sec:effondrement}

\begin{enclair}
Notre machine a une faille énorme. Si on la laisse faire, elle va trouver
une solution qui donne une erreur parfaitement nulle, sans rien apprendre du
tout. Il lui suffit de répondre toujours la même chose. C'est ce qu'on
appelle l'\emph{effondrement}, et c'est le cauchemar de tous ceux qui
entraînent ce genre de modèles.
\end{enclair}

Reprenons ce que nous demandons à la machine~: «~fais en sorte que ta
prédiction ressemble à la cible~». Et rappelons-nous du
chapitre~\ref{chap:apprendre}~: une machine qui apprend ne cherche pas à
bien faire. Elle cherche à faire baisser un nombre. Si elle trouve un moyen
malhonnête de faire baisser ce nombre, elle le prendra, sans le moindre
remords.

Or, il existe ici un moyen malhonnête, et il est même très facile à trouver.

\begin{definition}[Effondrement des représentations][effondrement]
Il y a \emph{effondrement} quand l'encodeur se met à sortir le même paquet
de nombres quelle que soit l'image qu'on lui donne. Les représentations ne
distinguent alors plus rien, et le modèle, bien que très content de lui,
est devenu complètement inutile.
\end{definition}

\begin{exemple}[Pourquoi la triche donne un zéro parfait]
Imaginons que les deux encodeurs se mettent d'accord pour toujours sortir le
paquet $(0 \,;\, 0 \,;\, 0)$, quoi qu'on leur montre. Et que le prédicteur,
de son côté, réponde toujours $(0 \,;\, 0 \,;\, 0)$ lui aussi.

Calculons l'erreur. On mesure l'écart en additionnant les différences prises
en valeur absolue, c'est-à-dire sans se soucier du signe~:

\[
  |0 - 0| \;+\; |0 - 0| \;+\; |0 - 0| \;=\; 0 .
\]

L'erreur vaut zéro. C'est le meilleur score possible, celui qu'aucun modèle
honnête n'atteindra jamais. Et pourtant ce modèle-là ne sait strictement
rien~: il donne la même réponse pour un chat, une rue et une plage.
\end{exemple}

\begin{schema}{L'effondrement. Trois vidéos qui n'ont rien à voir entre
elles ressortent de l'encodeur avec exactement le même paquet de nombres.
L'erreur est nulle, et le modèle est bon à jeter.}[jepa-effondrement]

\node[donnee, font=\sffamily\scriptsize, minimum width=15mm] (a) at (0,1.5)
     {un chat};
\node[donnee, font=\sffamily\scriptsize, minimum width=15mm] (b) at (0,0)
     {une rue};
\node[donnee, font=\sffamily\scriptsize, minimum width=15mm] (c) at (0,-1.5)
     {une plage};

\node[bloc, font=\sffamily\scriptsize, minimum width=21mm, minimum height=24mm]
     (e) at (3.5,0) {encodeur\\effondré};

\draw[fleche] (a) -- (e);
\draw[fleche] (b) -- (e);
\draw[fleche] (c) -- (e);

\node[bloc, font=\sffamily\scriptsize, draw=mbRose, fill=mbRose!8,
      minimum width=19mm] (s1) at (7.1,1.5) {$(0\,;0\,;0)$};
\node[bloc, font=\sffamily\scriptsize, draw=mbRose, fill=mbRose!8,
      minimum width=19mm] (s2) at (7.1,0)   {$(0\,;0\,;0)$};
\node[bloc, font=\sffamily\scriptsize, draw=mbRose, fill=mbRose!8,
      minimum width=19mm] (s3) at (7.1,-1.5){$(0\,;0\,;0)$};

\draw[fleche] (e) -- (s1);
\draw[fleche] (e) -- (s2);
\draw[fleche] (e) -- (s3);

\draw[accolade] (8.2,-1.9) -- (8.2,1.9);
\node[annotation, align=center, anchor=west] at (8.6,0)
     {toujours\\la même\\réponse};

\node[annotation, align=center] at (3.5,-2.4)
     {erreur nulle, savoir nul};
\end{schema}

\begin{intuition}
C'est l'élève qui a compris que le professeur note la ressemblance entre les
deux copies, et pas leur contenu. Il rend une feuille blanche. Le
professeur, qui est en réalité le même élève, rend lui aussi une feuille
blanche. Les deux copies se ressemblent parfaitement. Note maximale, et
aucun apprentissage.
\end{intuition}

Voilà pourquoi la question posée tout à l'heure était si bonne. Si l'on
utilisait \emph{un seul} encodeur des deux côtés, et qu'on le laissait
corriger ses boutons librement à partir des deux chemins, il découvrirait
cette triche en quelques centaines de pas. Nous le verrons se produire pour
de vrai, en direct, dans la section de code de ce chapitre.

Il faut donc empêcher la machine de tricher. Les auteurs emploient pour
cela deux garde-fous, et nous les prenons l'un après l'autre.

\section{Premier garde-fou : le gradient bloqué}

\begin{enclair}
Quand la machine se trompe, une correction remonte à l'envers dans tout le
montage pour tourner les boutons dans le bon sens. Le premier garde-fou
consiste à poser une barrière sur le chemin du professeur~: la correction a
le droit de remonter du côté de l'élève, mais pas du côté du corrigé. On
interdit ainsi à la machine de bouger la cible qu'elle est censée viser.
\end{enclair}

Rappelons d'abord d'où vient ce mot de gradient, vu au
chapitre~\ref{chap:apprendre}. Après avoir mesuré l'erreur, on calcule pour
chaque bouton dans quel sens et de combien il faudrait le tourner pour que
l'erreur diminue. Cette information voyage en sens inverse du calcul, de la
sortie vers l'entrée. C'est le gradient.

Maintenant, réfléchissons à ce que la machine ferait si on la laissait
libre. Elle doit rapprocher deux choses~: la prédiction et la cible. Il y a
deux façons de rapprocher deux points~:

\begin{enumerate}
  \item déplacer la prédiction vers la cible — c'est ce que l'on veut,
        puisque cela oblige l'élève à comprendre la scène~;
  \item déplacer la cible vers la prédiction — c'est la triche, puisqu'il
        suffit alors d'écraser tout le monde au même endroit.
\end{enumerate}

\begin{exemple}[Les deux façons de refermer un écart]
Prenons un seul nombre, pour voir la chose en face. La prédiction vaut
$0{,}5$ et la cible vaut $0{,}9$. L'écart entre les deux vaut donc
$0{,}9 - 0{,}5 = 0{,}4$.

\textbf{La façon honnête}~: on fait monter la prédiction de $0{,}5$ jusqu'à
$0{,}9$. L'écart tombe à $0$. Pour y arriver, l'élève a bien été obligé de
mieux regarder l'image.

\textbf{La façon tricheuse}~: on fait descendre la cible de $0{,}9$ jusqu'à
$0{,}5$. L'écart tombe à $0$ lui aussi, et pourtant personne n'a rien
appris. Pire~: si on laisse faire, la prédiction et la cible vont continuer
à se rapprocher l'une de l'autre jusqu'à valoir toutes les deux zéro.

Les deux façons donnent le même score. Voilà pourquoi il faut interdire la
seconde~: la machine, elle, ne voit aucune différence entre les deux.
\end{exemple}

La deuxième façon est de très loin la plus facile. La machine la
choisirait tout de suite. On la lui interdit donc, tout simplement.

\begin{definition}[Gradient bloqué, noté $\mathrm{sg}(\cdot)$][stop-gradient]
Le \emph{gradient bloqué} est une instruction qui dit~: «~calcule cette
valeur normalement, mais ne fais remonter aucune correction par ce
chemin~». On l'écrit $\mathrm{sg}(\cdot)$, d'après l'anglais
\emph{stop-gradient}, et l'on met entre les parenthèses ce que l'on veut
protéger. Dans notre montage, on protège la cible.
\end{definition}

\begin{schema}{Le trajet de la correction. Les flèches grises du haut sont
le calcul, qui va de gauche à droite. Les flèches orange du dessous sont la
correction, qui revient en arrière. La barre rouge lui interdit d'entrer
dans le chemin du professeur.}[jepa-stopgrad]

\node[bloc, font=\sffamily\scriptsize, minimum width=21mm] (e1) at (0,0)
     {encodeur élève};
\node[bloc accent, font=\sffamily\scriptsize, minimum width=21mm]
     (p1) at (3.7,0) {prédicteur};
\node[bloc, font=\sffamily\scriptsize, minimum width=16mm, draw=mbAmber,
      fill=mbAmber!8] (l1) at (7.4,0) {la perte};
\node[bloc calcul, font=\sffamily\scriptsize, minimum width=21mm]
     (e2) at (3.7,-3.0) {encodeur professeur};

% --- le calcul, vers la droite ---
\draw[fleche] (e1) -- (p1);
\draw[fleche] (p1) -- (l1);
\draw[fleche] (e2) -| (l1);
\node[etiquette] at (1.85,0.42) {le calcul avance};

% --- la correction, vers la gauche, en dessous ---
\draw[fleche accent, dashed] (l1.south) to[bend left=22] (p1.south);
\draw[fleche accent, dashed] (p1.south) to[bend left=22] (e1.south);
\node[annotation, align=center] at (3.7,-1.5)
     {la correction revient par ici};

% --- la barrière ---
\draw[line width=2pt, draw=mbRose!85!black] (6.6,-3.28) -- (6.6,-2.72);
\node[annotation, text=mbRose!85!black, align=center] at (6.6,-3.85)
     {barrière : la correction\\n'a pas le droit de passer};
\end{schema}

\begin{piege}
C'est l'erreur la plus fréquente, et surtout la plus silencieuse, de tout ce
domaine. Si vous oubliez le gradient bloqué, votre programme ne plantera
pas. Il n'affichera aucun message. La courbe d'erreur descendra même
magnifiquement vers zéro, ce qui vous rendra très fier. Et vous aurez
entraîné, pendant des jours, un modèle qui répond la même chose à tout.
Nous verrons au chapitre~25 comment repérer ce désastre à temps.
\end{piege}

Ce premier garde-fou est indispensable, mais il ne suffit pas tout à fait. Il
empêche la cible de reculer, mais il ne dit rien de la manière dont le
professeur, lui, doit évoluer. C'est le rôle du second garde-fou.

\section{Second garde-fou : le professeur qui suit de loin}
\label{sec:ema}

\begin{enclair}
Le professeur n'apprend pas tout seul~: il recopie l'élève. Mais il le
recopie très lentement, en n'attrapant à chaque pas qu'une toute petite
part du chemin. Résultat~: il est toujours en retard sur l'élève, et
c'est exactement ce retard qui empêche les deux de se mettre d'accord sur
une réponse vide, c'est-à-dire une réponse qui ne dit rien de ce que l'on
regarde.
\end{enclair}

La recette porte un nom un peu impressionnant, mais l'opération tient en une
seule ligne de calcul.

\begin{definition}[Moyenne mobile exponentielle (EMA)][ema]
La \emph{moyenne mobile exponentielle} est une façon de mettre à jour une
valeur en gardant beaucoup de l'ancienne et très peu de la nouvelle. On
choisit un nombre $m$ compris entre $0$ et $1$, appelé
\emph{coefficient}, et à chaque pas on remplace l'ancienne valeur par~:

\[
  \bar{\theta} \;\leftarrow\; m \times \bar{\theta} \;+\; (1 - m) \times \theta .
\]

La flèche $\leftarrow$ se lit «~devient~». À gauche du signe~$+$, ce que
l'on garde du professeur. À droite, ce que l'on prend de l'élève. Les trois
lettres EMA sont les premières lettres des trois mots anglais
\emph{Exponential Moving Average}, qui veulent dire exactement cela~:
moyenne mobile exponentielle.
\end{definition}

Ce nom fait peur, alors regardons-le de près~: il contient trois mots, et
chacun dit une chose simple. \emph{Moyenne}, parce que la nouvelle valeur
est un mélange de l'ancienne et de celle de l'élève. \emph{Mobile}, parce
que l'on refait ce mélange à chaque pas, sans jamais s'arrêter. Et
\emph{exponentielle}, parce qu'un vieux réglage ne disparaît pas d'un coup~:
il est multiplié par $m$ à chaque pas, donc il compte de moins en moins,
mais toujours un tout petit peu. Au bout de dix pas, ce qu'il en restait a
été multiplié dix fois de suite par $m$. Voilà tout ce que ce grand mot veut
dire.

Dans V-JEPA~2.1, ce coefficient vaut $m = 0{,}99925$. Traduisons ce nombre
en français courant~: à chaque pas, le professeur garde
$99{,}925\,\%$ de ce qu'il était, et n'attrape que $0{,}075\,\%$ de l'élève.
C'est très, très peu.

\begin{exemple}[Trois pas de moyenne mobile, calculés à la main]
Prenons un seul bouton, pour que tout tienne sur une ligne. Le professeur
est réglé sur $1{,}0$ et l'élève sur $2{,}0$. L'élève ne bouge pas, pour
simplifier. Faisons trois pas.

\textbf{Premier pas}~:
\[
  0{,}99925 \times 1{,}0 \;+\; 0{,}00075 \times 2{,}0
  \;=\; 0{,}99925 + 0{,}0015 \;=\; 1{,}0007500 .
\]

\textbf{Deuxième pas}~: on recommence, mais en partant de $1{,}00075$~:
\[
  0{,}99925 \times 1{,}0007500 \;+\; 0{,}0015 \;=\; 1{,}0014994 .
\]

\textbf{Troisième pas}~:
\[
  0{,}99925 \times 1{,}0014994 \;+\; 0{,}0015 \;=\; 1{,}0022483 .
\]

Regardons maintenant l'écart qui sépare le professeur de l'élève~:

\begin{center}
\footnotesize
\begin{tabular}{lrr}
\toprule
& \textbf{Professeur} & \textbf{Écart avec l'élève} \\
\midrule
Au départ      & $1{,}0000000$ & $1{,}0000000$ \\
Après 1 pas    & $1{,}0007500$ & $0{,}9992500$ \\
Après 2 pas    & $1{,}0014994$ & $0{,}9985006$ \\
Après 3 pas    & $1{,}0022483$ & $0{,}9977517$ \\
\bottomrule
\end{tabular}
\end{center}

En trois pas, l'écart n'a diminué que de deux millièmes. Autrement dit~: le
professeur n'a presque pas bougé. Et ce n'est pas un défaut, c'est
précisément ce que l'on cherche.
\end{exemple}

Pour se faire une idée du rythme, deux repères sont bien utiles.

\textbf{Repère un~: le temps de demi-parcours.} Il faut environ $924$ pas
pour que l'écart entre le professeur et l'élève soit divisé par deux. Dans
notre exemple, le professeur passe alors de $1{,}0$ à $1{,}5$, soit la
moitié du chemin vers $2{,}0$.

\textbf{Repère deux~: la mémoire.} On calcule $1 \div 0{,}00075 \approx
1\,333$. On dit que le professeur «~se souvient~» d'environ $1\,333$ pas
d'entraînement~: sa valeur est une sorte de moyenne des mille trois
cents derniers réglages de l'élève.

\begin{graphe}{Le professeur rattrape l'élève, mais très lentement.
L'élève est fixé à $2$ (trait orange). Le professeur part de $1$~: il lui
faut $924$ pas pour faire la moitié du chemin, et près de cinq mille pas
pour arriver tout près.}[jepa-ema]
\begin{axis}[mbaxis,
    xlabel={pas d'entraînement},
    ylabel={valeur du bouton},
    xmin=0, xmax=5000, ymin=0.9, ymax=2.15,
    legend pos=north west,
]
% repères du demi-parcours, tracés avant les courbes
\addplot[mbSlate!60, dashed, thin, forget plot]
        coordinates {(924,0.9) (924,1.5)};
\addplot[mbSlate!60, dashed, thin, forget plot]
        coordinates {(0,1.5) (924,1.5)};

\addplot[mbBlue, thick, domain=0:5000, samples=200, no marks]
        {2 - exp(-0.00075*x)};
\addlegendentry{le professeur}
\addplot[mbAmber, thick, dashed, domain=0:5000, samples=2, no marks] {2};
\addlegendentry{l'élève}

\node[circle, fill=mbRose, inner sep=1.6pt] at (axis cs:924,1.5) {};
\node[anchor=west, font=\sffamily\scriptsize, text=mbRose!85!black]
     at (axis cs:1100,1.42) {924 pas : la moitié du chemin};
\end{axis}
\end{graphe}

Reste à comprendre \emph{pourquoi} cette lenteur protège de l'effondrement.
L'explication tient en une phrase~: pour que la triche fonctionne, il
faudrait que l'élève et le professeur s'écrasent \emph{ensemble}, au même
moment, sur la même réponse vide. Mais le professeur ne suit qu'avec un
millier de pas de retard. Dès que l'élève se met à donner des réponses toutes pareilles, le
professeur, lui, continue de fournir des cibles variées, héritées d'un passé
où l'élève ne trichait pas encore. L'élève reçoit donc immédiatement une
erreur, et la triche ne lui rapporte plus rien.

\begin{piege}
Le coefficient $m$ se règle avec prudence, car les deux extrêmes sont
mauvais, et de façons opposées.

Si $m$ est trop petit — disons $0{,}5$ — le professeur rattrape l'élève
presque instantanément. Les deux deviennent le même appareil, et l'on
retombe droit dans l'effondrement.

Si $m$ est trop grand — disons $0{,}99999999$ — le professeur ne bouge
pratiquement plus. Il continue à donner les cibles de son premier jour, où
ses boutons étaient encore tirés au hasard. L'élève passe alors son temps à
imiter du bruit.
\end{piege}

\begin{aretenir}
Les deux garde-fous ne font pas la même chose, et il faut les deux. Le
gradient bloqué empêche la machine de \emph{tirer la cible vers elle}. La
moyenne mobile fait avancer la cible \emph{avec du retard}, pour que
personne ne puisse s'entendre sur une réponse vide.
\end{aretenir}

Nos deux garde-fous sont en place. Voyons maintenant à quoi ressemble un
pas d'entraînement complet, du début à la fin.

\section{Un pas d'entraînement, du début à la fin}

Voici l'enchaînement exact, dans l'ordre où vous l'écrirez dans votre
programme. Ces neuf étapes se répètent plus de cent mille fois. Nous les
retrouverons, presque mot pour mot, au chapitre~24, quand nous coderons pour
de bon.

\begin{enumerate}
  \item On prend un lot de vidéos.
  \item On tire un masque au hasard, et l'on fabrique la vue trouée $x$. La
        vue complète $y$, elle, est la vidéo telle quelle.
  \item L'encodeur élève regarde $x$ et produit les jetons de contexte.
  \item On ajoute les jetons de masque, qui portent les adresses des trous.
  \item Le prédicteur produit une réponse pour chaque position.
  \item L'encodeur professeur regarde $y$ et produit les cibles. On pose ici
        le gradient bloqué.
  \item On mesure l'écart entre les réponses et les cibles~: c'est la perte.
  \item La correction remonte, et l'on tourne les boutons de l'élève et du
        prédicteur.
  \item On met à jour le professeur avec la moyenne mobile. Puis on
        recommence au point~1.
\end{enumerate}

Ces neuf étapes forment une boucle~: la neuvième ramène à la première, et
l'on repart pour un tour. La figure~\ref{fig:jepa-boucle} la montre d'un
seul coup d'œil, en regroupant les neuf étapes en quatre moments.

\begin{schema}{La boucle d'entraînement. Les neuf étapes se rangent en
quatre moments, et le quatrième ramène au premier. C'est ce tour complet
que la machine répète plus de cent mille fois.}[jepa-boucle]

\node[bloc, font=\sffamily\scriptsize, minimum width=30mm] (a) at (0,0)
     {\textbf{1--2. Préparer}\\on cache des morceaux\\et l'on obtient $x$ et $y$};
\node[bloc accent, font=\sffamily\scriptsize, minimum width=30mm]
     (b) at (5.4,0)
     {\textbf{3--5. Deviner}\\l'élève regarde, le\\prédicteur répond};
\node[bloc, font=\sffamily\scriptsize, minimum width=30mm, draw=mbAmber,
      fill=mbAmber!8] (c) at (5.4,-2.6)
     {\textbf{6--7. Comparer}\\on mesure l'écart\\avec les cibles};
\node[bloc calcul, font=\sffamily\scriptsize, minimum width=30mm]
     (d) at (0,-2.6)
     {\textbf{8--9. Corriger}\\on tourne les boutons,\\puis le professeur suit};

\draw[fleche] (a) -- (b);
\draw[fleche] (b) -- (c);
\draw[fleche] (c) -- (d);
\draw[fleche accent] (d) -- (a);

\node[annotation, align=center] at (2.7,-1.3)
     {et l'on recommence\\encore et encore};
\end{schema}

\begin{surlesdonnees}
Pour donner l'échelle réelle~: dans V-JEPA~2.1, cette boucle tourne
$135\,000$ fois pendant la phase principale, puis $12\,000$ fois de plus
pendant une phase finale à plus haute résolution. Cela fait $147\,000$
répétitions au total, sur des centaines de machines à la fois. Avec un
coefficient de $0{,}99925$ et un temps de demi-parcours de $924$ pas, l'écart
entre le professeur et l'élève aura eu le temps d'être divisé par deux
environ cent soixante fois, puisque $147\,000 \div 924 \approx 159$. Le
professeur est donc très lent à l'échelle d'un pas, mais pas du tout à
l'échelle de l'entraînement complet.
\end{surlesdonnees}

Assez de théorie. Faisons tourner tout cela.

\section{Passons au code}

Nous allons écrire trois petits programmes. Le premier refait le calcul de
la moyenne mobile, pour vérifier les trois nombres que nous venons de poser
à la main. Le deuxième
monte un entraînement \emph{sans} garde-fou, et nous regardons le modèle
s'effondrer sous nos yeux. Le troisième remet les deux garde-fous, et nous
vérifions que tout tient. Tout cela tourne en quelques secondes sur
n'importe quel ordinateur.

Commençons par la moyenne mobile.

\begin{code}{python}[La moyenne mobile du professeur, pas à pas]
m = 0.99925
professeur, eleve = 1.0, 2.0

for pas in range(1, 4):
    professeur = m * professeur + (1 - m) * eleve
    print(f"pas {pas} : professeur = {professeur:.7f}"
          f"   ecart = {eleve - professeur:.7f}")
\end{code}

À l'exécution, on obtient exactement les nombres du tableau de tout à
l'heure~:

\begin{sortie}
pas 1 : professeur = 1.0007500   ecart = 0.9992500
pas 2 : professeur = 1.0014994   ecart = 0.9985006
pas 3 : professeur = 1.0022483   ecart = 0.9977517
\end{sortie}

Si l'on laisse tourner la boucle plus longtemps, on retrouve aussi nos deux
repères~: au pas $924$, le professeur vaut $1{,}5001$, soit presque exactement
la moitié du chemin~; au pas $5\,000$, il vaut $1{,}9765$.

Passons maintenant à l'effondrement. Pour le voir arriver, nous avons besoin
d'un instrument de mesure, et c'est important~: la perte, à elle seule, ne
nous préviendra de rien, puisqu'elle descendra joliment. Nous allons donc
mesurer autre chose, que nous appellerons la \emph{diversité}~: à quel point
les sorties de l'encodeur diffèrent les unes des autres. Si ce nombre tombe
vers zéro, c'est que toutes les réponses se ressemblent, et donc que le
modèle s'est effondré.

Cette mesure repose sur une petite opération que nous n'avons pas encore
rencontrée, et qu'il faut donc expliquer avant de s'en servir.

\begin{definition}[L'écart-type][ecart-type]
L'\emph{écart-type} d'une liste de nombres dit à quel point ces nombres sont
étalés. On regarde d'abord leur moyenne, puis de combien chacun s'en écarte.
Si tous les nombres de la liste sont égaux, l'écart-type vaut exactement
zéro. Plus ils sont éparpillés, plus il est grand. En Python, on l'obtient
avec \co{python}{.std()}.
\end{definition}

\begin{exemple}[Deux listes, deux écarts-types]
La liste $(5 \,;\, 5 \,;\, 5)$ a pour moyenne $5$, et aucun de ses nombres
ne s'écarte de cette moyenne. Son écart-type vaut donc $0$. C'est exactement
la situation de l'effondrement~: tout le monde répond la même chose.

La liste $(3 \,;\, 5 \,;\, 7)$ a elle aussi pour moyenne $5$, mais ses
nombres s'en écartent de $2$, de $0$ et de $2$. Son écart-type vaut $2$, et
c'est bien ce que rend Python si vous le lui demandez. C'est la situation
d'un modèle en bonne santé~: les réponses sont variées.
\end{exemple}

Notre diversité n'est rien d'autre que cela~: l'écart-type des sorties de
l'encodeur, mesuré d'un exemple à l'autre, puis moyenné pour n'avoir qu'un
seul nombre à surveiller.

\begin{code}{python}[Un entraînement sans le moindre garde-fou]
import torch
import torch.nn as nn

torch.manual_seed(0)

encodeur = nn.Sequential(nn.Linear(8, 16), nn.GELU(), nn.Linear(16, 16))
predicteur = nn.Sequential(nn.Linear(16, 16), nn.GELU(), nn.Linear(16, 16))
optim = torch.optim.Adam(
    list(encodeur.parameters()) + list(predicteur.parameters()), lr=1e-3)

x = torch.randn(64, 8)             # a la place de la vue trouee : du hasard
y = x + 0.1 * torch.randn(64, 8)   # a la place de la vue complete

for pas in range(1, 301):
    cible = encodeur(y)                # AUCUN gradient bloque ici
    prediction = predicteur(encodeur(x))
    perte = (prediction - cible).abs().mean()

    optim.zero_grad()
    perte.backward()
    optim.step()

    if pas % 100 == 0 or pas == 1:
        diversite = cible.std(dim=0).mean().item()
        print(f"pas {pas:4d} : perte = {perte.item():.6f}"
              f"   diversite = {diversite:.6f}")
\end{code}

Voici ce que cela donne~:

\begin{sortie}[Le modèle s'effondre]
pas    1 : perte = 0.273260   diversite = 0.219396
pas  100 : perte = 0.039338   diversite = 0.066842
pas  200 : perte = 0.007948   diversite = 0.017721
pas  300 : perte = 0.002699   diversite = 0.007204
\end{sortie}

Lisons ces quatre lignes ensemble, car elles racontent toute l'histoire de
ce chapitre. La colonne de la perte descend vers zéro~: du point de vue de
la machine, tout va très bien, elle fait exactement ce qu'on lui a demandé.
Mais celle de la diversité descend elle aussi, et c'est la catastrophe~:
partie de
$0{,}22$, la diversité est tombée à $0{,}007$, soit trente fois moins. Les
sorties se sont écrasées les unes sur les autres. En trois cents pas à
peine, sur un ordinateur ordinaire, nous venons de reproduire le désastre
que nous décrivions plus haut.

Remettons maintenant les deux garde-fous. Trois choses changent, et
seulement trois~: un professeur séparé, dont les boutons ne sont jamais
tournés par la correction~; un gradient bloqué autour de la cible~; et la mise à
jour par moyenne mobile à la fin de chaque pas.

\begin{code}{python}[Le même entraînement, avec les deux garde-fous]
eleve = nn.Sequential(nn.Linear(8, 16), nn.GELU(), nn.Linear(16, 16))
predicteur = nn.Sequential(nn.Linear(16, 16), nn.GELU(), nn.Linear(16, 16))

professeur = nn.Sequential(nn.Linear(8, 16), nn.GELU(), nn.Linear(16, 16))
professeur.load_state_dict(eleve.state_dict())   # il part de l'eleve
for p in professeur.parameters():
    p.requires_grad_(False)                      # il n'apprendra jamais seul

optim = torch.optim.Adam(
    list(eleve.parameters()) + list(predicteur.parameters()), lr=1e-3)
m = 0.996                                        # coefficient de la moyenne

x = torch.randn(64, 8)                           # de nouvelles donnees
y = x + 0.1 * torch.randn(64, 8)

for pas in range(1, 301):
    with torch.no_grad():                        # <-- le gradient bloque
        cible = professeur(y)

    prediction = predicteur(eleve(x))
    perte = (prediction - cible).abs().mean()

    optim.zero_grad()
    perte.backward()
    optim.step()

    with torch.no_grad():                        # <-- la moyenne mobile
        for pp, pe in zip(professeur.parameters(), eleve.parameters()):
            pp.mul_(m).add_(pe, alpha=1 - m)

    if pas % 100 == 0 or pas == 1:
        diversite = cible.std(dim=0).mean().item()
        print(f"pas {pas:4d} : perte = {perte.item():.6f}"
              f"   diversite = {diversite:.6f}")
\end{code}

Et le résultat~:

\begin{sortie}[Le modèle tient debout]
pas    1 : perte = 0.247933   diversite = 0.171066
pas  100 : perte = 0.069328   diversite = 0.188252
pas  200 : perte = 0.048201   diversite = 0.232273
pas  300 : perte = 0.038985   diversite = 0.275197
\end{sortie}

Comparons les deux tableaux, colonne par colonne. La perte descend dans les
deux cas — c'est bien la preuve qu'elle ne nous apprend rien à elle seule.
Mais la diversité, elle, fait le contraire~: elle tombait à $0{,}007$ sans
garde-fou, elle \emph{monte} à $0{,}275$ avec. Les représentations, au lieu
de s'écraser, se sont écartées les unes des autres. Le modèle a commencé à
distinguer ses entrées, et c'est exactement le comportement que nous
voulions.

\begin{plusloin}
Vous avez peut-être remarqué que nous avons pris $m = 0{,}996$ dans ce petit
programme, et non $0{,}99925$ comme dans l'article. C'est volontaire, et
c'est une règle générale~: le bon coefficient dépend du nombre de pas que
dure l'entraînement. Ici nous n'en faisons que trois cents, alors que
l'article en fait cent quarante-sept mille. Avec $0{,}99925$, notre
professeur n'aurait pas eu le temps de bouger. Essayez, vous verrez~: la
diversité ne monte alors que de $0{,}171$ à $0{,}189$, au lieu de $0{,}275$.
Elle ne s'effondre pas, mais elle n'avance presque plus, faute de temps.
\end{plusloin}

\begin{aretenir}
Une architecture JEPA tient en cinq pièces et deux garde-fous. Les pièces~:
la vue trouée $x$, la vue complète $y$, l'encodeur élève $E_{\theta}$,
l'encodeur professeur $E_{\bar{\theta}}$ et le prédicteur $P_{\varphi}$. Les
garde-fous~: le gradient bloqué, qui empêche de tirer la cible vers soi, et
la moyenne mobile, qui fait avancer le professeur avec du retard. Et une
seule idée au-dessus de tout cela~: on compare des représentations, jamais
des pixels.
\end{aretenir}

\vspace{0.5em}

Nous savons désormais comment la machine est bâtie, et pourquoi elle ne
s'écroule pas. Il nous reste à écrire, noir sur blanc, ce que
l'entraînement cherche à faire baisser. Cela tient en une seule ligne de
symboles~: c'est la formule officielle de V-JEPA~2. Nous la découperons
symbole par symbole, sans en laisser un seul de côté, au chapitre suivant. Nous y verrons non seulement
\emph{quoi} la machine doit prédire, mais aussi, et c'est le point capital,
\emph{à quels endroits} on lui demande des comptes. Car un détail se cache
là — un tout petit détail, presque invisible — et c'est de ce détail que
naîtra V-JEPA~2.1.

\phantomsection
\section*{Exercices}
\addcontentsline{toc}{section}{Exercices}

\begin{exercice}[Deux pas de moyenne mobile]
Un professeur est réglé sur $3{,}0$, un élève sur $5{,}0$, et le coefficient
vaut $m = 0{,}99$. Calculez la valeur du professeur après un pas, puis après
deux pas. De combien l'écart a-t-il diminué~?
\end{exercice}

\begin{solution}
Premier pas~: $0{,}99 \times 3{,}0 + 0{,}01 \times 5{,}0 = 2{,}97 + 0{,}05
= 3{,}02$.

Deuxième pas~: $0{,}99 \times 3{,}02 + 0{,}05 = 2{,}9898 + 0{,}05 =
3{,}0398$.

L'écart valait $5{,}0 - 3{,}0 = 2{,}0$ au départ. Il vaut maintenant
$5{,}0 - 3{,}0398 = 1{,}9602$. Il a donc diminué de $0{,}0398$, soit environ
deux pour cent en deux pas. Remarquez au passage que $m = 0{,}99$ est
beaucoup plus rapide que le $0{,}99925$ de l'article~: le professeur y
attrape un centième de l'élève à chaque pas, contre moins d'un millième dans
V-JEPA~2.1.
\end{solution}

\begin{exercice}[Compter les jetons du prédicteur]
Un clip contient $4\,096$ jetons. Le masque en recouvre $2\,500$. Combien de
jetons l'encodeur élève reçoit-il~? Combien le prédicteur reçoit-il en
entrée~? Combien en ressort-il~?
\end{exercice}

\begin{solution}
L'encodeur élève ne voit que le contexte, soit $4\,096 - 2\,500 = 1\,596$
jetons.

Le prédicteur reçoit ces $1\,596$ jetons de contexte, plus $2\,500$ jetons
de masque, soit $1\,596 + 2\,500 = 4\,096$ jetons en entrée.

Et il en ressort $4\,096$, puisqu'il produit une réponse par jeton reçu.
Gardez bien en tête que $1\,596$ de ces réponses concernent des endroits qui
étaient déjà visibles~: nous verrons bientôt que V-JEPA~2 les jette, et que
c'est là son point faible.
\end{solution}

\begin{exercice}[Trouver l'erreur]
Un collègue vous montre sa boucle d'entraînement. Elle contient la ligne
suivante~:

\begin{center}
\term{cible = professeur(y)}
\end{center}

Sa courbe d'erreur descend magnifiquement vers zéro en deux heures, et il
est très content. Que lui dites-vous~?
\end{exercice}

\begin{solution}
Qu'il a très probablement oublié le gradient bloqué, et qu'il faut écrire
cette ligne à l'intérieur d'un \term{with torch.no\_grad():}.

Il faut aussi lui expliquer pourquoi sa jolie courbe est un mauvais signe et
non un bon~: une erreur qui descend vers zéro en deux heures, sur un
problème difficile, ne veut pas dire que le modèle a compris quelque chose.
Cela veut dire, presque toujours, qu'il a trouvé une réponse vide.
Conseillez-lui de tracer la diversité de ses représentations à côté de sa
perte~: c'est ce deuxième nombre qui dit la vérité.
\end{solution}

\begin{exercice}[Le compte des nombres]
Un lecteur affirme~: «~JEPA est meilleur parce qu'il a moins de nombres à
deviner qu'une méthode qui redessine les pixels~». A-t-il raison~?
\end{exercice}

\begin{solution}
Non, et c'est le piège classique de ce chapitre. Rien ne garantit qu'une
représentation contienne moins de nombres qu'une zone de pixels~; elle peut
même en contenir beaucoup plus.

Ce qui change, ce n'est pas la quantité, c'est la nature de ce qu'on
demande. Les pixels sont imposés du dehors, avec tout leur hasard~: la
position exacte de chaque brin d'herbe doit être devinée, alors qu'elle est
indevinable. La représentation, elle, est fabriquée par le modèle lui-même,
qui a donc le droit d'y laisser tomber ce qui n'est pas prévisible et de n'y
garder que ce qui compte.
\end{solution}

\begin{exercice}[Régler le coefficient]
Vous choisissez $m = 0{,}5$ pour la moyenne mobile de votre professeur. Que
va-t-il se passer, et pourquoi~?
\end{exercice}

\begin{solution}
Avec $m = 0{,}5$, le professeur attrape la moitié de l'élève à chaque pas.
En cinq ou six pas, il lui est devenu presque identique.

Or, c'est justement le retard du professeur qui empêchait l'effondrement.
Sans ce retard, l'élève et le professeur peuvent s'écraser ensemble sur une
réponse vide, et plus rien ne les en empêche. Vous verrez donc la perte
plonger vers zéro et la diversité des représentations s'effondrer, exactement
comme dans notre deuxième programme.
\end{solution}

\begin{exercice}[À quoi servent les adresses]
Que se passerait-il si les jetons de masque ne portaient pas l'adresse du
trou, mais étaient tous rigoureusement identiques~?
\end{exercice}

\begin{solution}
Le prédicteur recevrait plusieurs enveloppes vides toutes
pareilles, et il n'aurait aucun moyen de savoir laquelle correspond
au coin en haut à gauche et laquelle correspond au bas de l'image.

N'ayant pas de raison de les traiter différemment, il donnerait à peu près
la même réponse pour tous les trous — sans doute une sorte de moyenne de ce
qu'il y a d'habitude dans une vidéo. La perte ne descendrait pas très bas, et
surtout le modèle n'apprendrait rien de la façon dont la scène est rangée
dans l'espace.
L'adresse n'est donc pas un détail de rangement~: c'est elle qui rend la
question possible à poser.
\end{solution}
